	\section{Análisis y Diseño}
	\label{cap:analisis_diseno}
	Este capítulo documenta el análisis y el diseño de la plataforma tal como quedó definida al
	cierre de esta etapa del Trabajo Terminal. Su propósito es dejar constancia de \emph{qué} se
	construyó y, sobre todo, de \emph{por qué} se construyó de esa manera: los requerimientos que
	se levantaron con la Defensoría, los actores que intervienen, los casos de uso que delimitan
	el alcance funcional, la arquitectura que soporta la solución y el modelo de datos que la
	sustenta. La construcción concreta de cada componente ---el código, el despliegue y las
	pruebas--- se documenta en el Capítulo~\ref{cap:implementacion}, y el recorrido funcional de
	cada módulo entregado en el Capítulo~\ref{cap:modulos}.
	
	Los artefactos que aquí se presentan siguen el orden en que fueron producidos: primero el
	levantamiento de requerimientos con el cliente, del que se derivaron los casos de uso; después
	el modelo de estados del expediente, que ordena el flujo de trabajo; sobre esa base la
	arquitectura de componentes y el modelo de datos; y por último las decisiones de diseño con
	las alternativas que se descartaron en cada punto. Este orden no es una convención editorial:
	cada artefacto es insumo del siguiente, y esa dependencia es la que sostiene la coherencia
	técnica del sistema.
	
	\subsection{Descripción general de la solución}
	\label{sec:descripcion_solucion}
	La solución es una plataforma web para la gestión digital de quejas y denuncias ante la
	Defensoría de los Derechos Politécnicos. Permite a cualquier integrante de la comunidad
	politécnica presentar una queja y darle seguimiento ---con o sin cuenta de usuario--- y provee
	al personal de la Defensoría las herramientas para recibirla, validarla, canalizarla y
	administrar la operación técnica del sistema.
	
	Al cierre de esta etapa la plataforma está compuesta por siete microservicios independientes
	desarrollados con Spring Boot que comparten una única base de datos PostgreSQL, tres
	aplicaciones web independientes desarrolladas con Angular ---una por tipo de usuario--- y
	cuatro contenedores de Nginx que actúan como proxy inverso público y como servidores de
	archivos estáticos. Todo el conjunto se ejecuta en contenedores sobre dos servidores virtuales
	privados, con acceso público cifrado mediante TLS y certificado renovado de forma automática.
	
	\subsection{Actores del sistema}
	\label{sec:actores}
	El análisis identificó tres actores con necesidades, permisos e interfaces claramente
	diferenciadas. Esta separación es la que después justifica que existan tres aplicaciones web
	distintas en lugar de una sola con vistas condicionadas por rol: cada actor tiene un perfil de
	uso, un nivel de confianza y una superficie de exposición distintos, y mantenerlos en
	aplicaciones separadas reduce el riesgo de que un defecto en una interfaz exponga
	funcionalidad de otra.
	
	\begin{itemize}
		\item \textbf{Quejoso.} Cualquier integrante de la comunidad politécnica ---estudiante,
		personal académico o administrativo--- que presenta una queja o le da seguimiento. Es el
		único actor externo y el único que puede operar sin credenciales previas.
	
		\item \textbf{Personal de recepción.} Personal de la Defensoría encargado de recibir,
		validar y canalizar las quejas que llegan por el portal público o en papel.
	
		\item \textbf{Administrador de tecnologías de la información.} Responsable de la operación
		técnica y administrativa de la plataforma: cuentas del resto del personal, catálogo
		institucional, plantillas de documentos oficiales, respaldos de la base de datos y bitácora
		de seguridad.
	\end{itemize}
	
	El sistema define además cinco roles de personal, de los cuales solo dos cuentan con interfaz
	propia al cierre de esta etapa. La Tabla~\ref{tab:actores_roles} resume la correspondencia
	entre actores, roles e interfaces.
	
	\rowcolors{2}{rowblue}{white}
	
	\small
	\renewcommand{\arraystretch}{1.25}
	\setlength{\tabcolsep}{4pt}
	
	\begin{longtable}{|p{2.7cm}|p{4.3cm}|p{3.8cm}|p{3.8cm}|}
		\caption{Actores, roles e interfaces del sistema}
		\label{tab:actores_roles} \\
		\hline
		\rowcolor{headerblue}
		\headercell{Actor} & \headercell{Rol en el sistema} & \headercell{Responsabilidad} &
		\headercell{Interfaz} \\
		\hline
		\endfirsthead
	
		\hline
		\rowcolor{headerblue}
		\headercell{Actor} & \headercell{Rol en el sistema} & \headercell{Responsabilidad} &
		\headercell{Interfaz} \\
		\hline
		\endhead
	
		Quejoso & ---\ (usuario externo) & Presentar quejas y darles seguimiento &
		Portal público del quejoso \\ \hline
	
		Personal de recepción & \texttt{RECEPCIONISTA} & Recibir, validar, rechazar y turnar quejas &
		Panel de revisión \\ \hline
	
		Administrador de TI & \texttt{ADMIN\_SISTEMAS} & Operación técnica, cuentas, catálogos,
		respaldos y auditoría & Consola de administración \\ \hline
	
		Analista de primer contacto & \texttt{ANALISTA\_\allowbreak PRIMER\_\allowbreak CONTACTO} & Evaluación jurídica de
		competencia & Sin interfaz propia en esta etapa \\ \hline
	
		Subdefensor & \texttt{SUBDEFENSOR} & Solicitud de investigación y conciliación &
		Sin interfaz propia en esta etapa \\ \hline
	
		Defensor & \texttt{DEFENSOR} & Resolución y cierre del expediente &
		Sin interfaz propia en esta etapa \\ \hline
	\end{longtable}
	
	\normalsize
	\setlength{\tabcolsep}{6pt}
	\renewcommand{\arraystretch}{1}
	
	Los tres roles sin interfaz propia no son un descuido del diseño: existen como cuentas
	capturables desde la consola de administración y el sistema los emplea como catálogo de
	personal disponible al momento de turnar una queja. Lo que no tienen todavía es un panel con
	permisos propios, y esa limitación se declara de forma explícita en el alcance de esta etapa
	(Sección~\ref{sec:limitaciones}).
	
	\subsection{Especificación de requerimientos}
	\label{sec:requerimientos}
	El levantamiento de requerimientos no se limitó a la lectura de la documentación oficial de la
	Defensoría, porque la operación real de una organización rara vez coincide por completo con lo
	que sus documentos formales describen. El equipo sostuvo una serie de reuniones con el personal
	de la institución que evolucionaron de forma iterativa: cada sesión aclaró aspectos que la
	anterior había dejado abiertos, y cada avance técnico intermedio se convirtió en insumo de la
	siguiente reunión.
	
	La primera sesión tuvo como objetivo comprender el contexto operativo e identificar las áreas
	que intervienen en la gestión de una queja. Con ese contexto el equipo elaboró un diagrama de
	proceso en notación BPMN a partir de los documentos oficiales, que se presentó en una segunda
	sesión como base de discusión. El resultado fue el hallazgo más valioso de todo el
	levantamiento: el personal de la Defensoría señaló que varias etapas del procedimiento formal
	no se ejecutaban en la práctica, de modo que se depuraron en conjunto hasta obtener un diagrama
	que refleja el flujo operativo real y no el prescrito. Validado ese diagrama, se establecieron
	sesiones por área para traducir las actividades de cada rol en casos de uso y vistas concretas,
	con sesiones de revisión intermedias para confirmar que los artefactos reflejaran lo acordado.
	
	El resultado acumulado se organizó en tres categorías ---requerimientos funcionales, no
	funcionales y reglas de negocio--- y se estructuró por actor, de modo que cada requerimiento
	tiene un dueño identificable. La nomenclatura empleada codifica el actor en el identificador:
	\texttt{RF-Q} para los requerimientos funcionales del quejoso, \texttt{RF-R} para los del
	personal de recepción y \texttt{RF-A} para los de administración, y de forma análoga
	\texttt{RNF-} y \texttt{RN-} para los no funcionales y las reglas de negocio. Esta convención
	permite rastrear cualquier elemento del diseño hasta el actor que lo motivó.
	
	\subsubsection{Requerimientos funcionales del quejoso}
	Los requerimientos del quejoso se agrupan en cinco bloques: operación sin cuenta, gestión de la
	cuenta, gestión de las quejas propias, seguimiento y servicios de apoyo. El primer bloque es el
	más determinante para el diseño, porque exige que el sistema acepte quejas de personas que no
	tienen ---ni necesariamente quieren tener--- una cuenta, lo que obliga a un punto de acceso
	público sin autenticación y a un mecanismo de identificación posterior basado en el folio.
	
	\small
	\renewcommand{\arraystretch}{1.2}
	\setlength{\tabcolsep}{4pt}
	
	\begin{longtable}{|p{1.5cm}|p{3.6cm}|p{9.0cm}|}
		\caption{Requerimientos funcionales del actor Quejoso}
		\label{tab:rf_quejoso} \\
		\hline
		\rowcolor{headerblue}
		\headercell{ID} & \headercell{Requerimiento} & \headercell{Descripción} \\
		\hline
		\endfirsthead
	
		\hline
		\rowcolor{headerblue}
		\headercell{ID} & \headercell{Requerimiento} & \headercell{Descripción} \\
		\hline
		\endhead
	
		RF-Q01 & Consultar una queja sin cuenta & Dado un folio y el correo con que se presentó,
		mostrar el detalle de la queja: motivo, descripción, estatus y fecha. \\ \hline
	
		RF-Q02 & Registrar una queja sin cuenta & Capturar la identidad completa del quejoso,
		unidad académica, fecha de los hechos, datos del denunciado, descripción libre y varios
		archivos de evidencia en una sola petición; generar un folio de seguimiento. \\ \hline
	
		RF-Q03 & Capturar datos del tutor & Si la fecha de nacimiento indica que el quejoso es menor
		de edad, exigir los datos de un tutor o adulto responsable antes de permitir el envío.
		\\ \hline
	
		RF-Q04 & Activar una cuenta de seguimiento & A partir de un folio y correo ya registrados en
		una queja existente, definir una contraseña y activar la cuenta. \\ \hline
	
		RF-Q05 & Iniciar sesión & Autenticarse con correo institucional y contraseña, obteniendo un
		token de sesión con vigencia limitada. \\ \hline
	
		RF-Q06 & Recuperar la contraseña & Restablecer la contraseña mediante un código de
		verificación de seis dígitos enviado por correo electrónico. \\ \hline
	
		RF-Q07 & Consultar un resumen de actividad & Mostrar al entrar al panel el conteo de quejas
		propias agrupadas por estatus. \\ \hline
	
		RF-Q08 & Listar y consultar las quejas propias & Listar las quejas del usuario con filtros
		por estatus y texto, y mostrar el detalle completo de cada una con su línea de tiempo de
		estatus. \\ \hline
	
		RF-Q09 & Editar una queja propia & Modificar descripción, unidad académica, fecha de hechos
		y datos del denunciado, únicamente mientras la queja no haya iniciado su validación.
		\\ \hline
	
		RF-Q10 & Registrar una queja autenticado & Presentar una queja nueva desde el panel sin
		volver a capturar la identidad, que se toma del token de sesión. \\ \hline
	
		RF-Q11 & Consultar las evidencias propias & Previsualizar y descargar los archivos adjuntos
		a una queja propia. \\ \hline
	
		RF-Q12 & Responder acuerdos de conciliación & Consultar los acuerdos propuestos por la
		Defensoría y aceptarlos o rechazarlos, con comentario opcional. \\ \hline
	
		RF-Q13 & Consultar el centro de notificaciones & Consultar un historial persistente de
		avisos con contador de no leídas y marcado individual como leída. \\ \hline
	
		RF-Q14 & Consultar y editar el perfil & Editar correo personal, teléfono, unidad académica y
		domicilio; consultar en modo de solo lectura el nombre, el correo institucional y la boleta o
		número de empleado. \\ \hline
	
		RF-Q15 & Consultar el catálogo institucional & Seleccionar la dependencia o unidad académica
		involucrada desde el catálogo del Instituto, en lugar de capturarla como texto libre.
		\\ \hline
	
		RF-Q16 & Consultar preguntas frecuentes & Consultar sin sesión un menú fijo de preguntas y
		respuestas sobre la Defensoría y el proceso de queja. \\ \hline
	
		RF-Q17 & Recibir avisos por correo & Recibir confirmación de creación de cuenta, confirmación
		de restablecimiento de contraseña y, en el caso de menores de edad, aviso al tutor.
		\\ \hline
	\end{longtable}
	
	\normalsize
	
	\subsubsection{Requerimientos funcionales del personal de recepción}
	El diseño de este bloque partió de una restricción que el propio personal planteó en las
	sesiones de levantamiento: no todas las quejas llegan por el portal. Una parte del volumen
	ingresa en papel, y el sistema debía poder incorporarlas al mismo flujo digital sin obligar al
	quejoso a capturarlas de nuevo. De ahí el requerimiento de registro manual (RF-R08), que es la
	razón por la que el modelo de datos distingue el origen de cada queja.
	
	\small
	\renewcommand{\arraystretch}{1.2}
	
	\begin{longtable}{|p{1.5cm}|p{3.6cm}|p{9.0cm}|}
		\caption{Requerimientos funcionales del actor Personal de recepción}
		\label{tab:rf_recepcion} \\
		\hline
		\rowcolor{headerblue}
		\headercell{ID} & \headercell{Requerimiento} & \headercell{Descripción} \\
		\hline
		\endfirsthead
	
		\hline
		\rowcolor{headerblue}
		\headercell{ID} & \headercell{Requerimiento} & \headercell{Descripción} \\
		\hline
		\endhead
	
		RF-R01 & Iniciar sesión & Autenticarse con las credenciales de personal administrativo desde
		el panel de revisión. \\ \hline
	
		RF-R02 & Consultar la bandeja de entrada & Consultar la lista de quejas por trabajar, con
		contadores de pendientes, en proceso y turnadas en el día. \\ \hline
	
		RF-R03 & Consultar el detalle de una queja & Consultar el resumen completo del expediente y
		los documentos adjuntos, para su validación. \\ \hline
	
		RF-R04 & Buscar antecedentes & Consultar otras quejas previas de la misma persona, para
		detectar posibles duplicados antes de turnar. \\ \hline
	
		RF-R05 & Consultar los documentos adjuntos & Previsualizar y descargar las evidencias
		asociadas a una queja. \\ \hline
	
		RF-R06 & Rechazar una queja & Rechazar una queja por documentación incompleta, indicando
		motivos y observaciones, con notificación automática al quejoso. \\ \hline
	
		RF-R07 & Turnar una queja & Canalizar la queja a un área o dependencia y a un defensor o
		subdefensor específico, con comentarios, generando el folio oficial de turnado. \\ \hline
	
		RF-R08 & Registrar una queja recibida en papel & Dar de alta manualmente una queja física
		con la identidad del quejoso, dependencia, número de oficio, fecha de recepción, tipo de
		documento, descripción, ubicación física del expediente y un archivo digitalizado opcional.
		\\ \hline
	
		RF-R09 & Consultar el historial de trámites & Consultar los trámites ya procesados, con
		filtros por texto, estatus y fecha. \\ \hline
	
		RF-R10 & Exportar el historial & Exportar el historial filtrado a un archivo de hoja de
		cálculo. \\ \hline
	
		RF-R11 & Emitir un acuerdo de conciliación & Emitir un acuerdo dirigido al quejoso de una
		queja, con asunto y términos. \\ \hline
	
		RF-R12 & Consultar catálogos de apoyo & Consultar la lista de dependencias y la de personal
		disponible para turnado. \\ \hline
	
		RF-R13 & Cambiar la propia contraseña & Modificar la contraseña de la cuenta propia desde el
		panel. \\ \hline
	\end{longtable}
	
	\normalsize
	
	\subsubsection{Requerimientos funcionales del administrador de TI}
	Este bloque concentra las funciones que hacen operable la plataforma a largo plazo sin
	intervención del equipo de desarrollo. Su requerimiento más sensible es el de respaldo y
	restauración (RF-A06), porque incluye la única operación destructiva de todo el sistema y por
	ello recibe un tratamiento de seguridad diferenciado.
	
	\small
	\renewcommand{\arraystretch}{1.2}
	
	\begin{longtable}{|p{1.5cm}|p{3.6cm}|p{9.0cm}|}
		\caption{Requerimientos funcionales del actor Administrador de TI}
		\label{tab:rf_admin} \\
		\hline
		\rowcolor{headerblue}
		\headercell{ID} & \headercell{Requerimiento} & \headercell{Descripción} \\
		\hline
		\endfirsthead
	
		\hline
		\rowcolor{headerblue}
		\headercell{ID} & \headercell{Requerimiento} & \headercell{Descripción} \\
		\hline
		\endhead
	
		RF-A01 & Iniciar sesión & Autenticarse como personal administrativo. \\ \hline
	
		RF-A02 & Consultar el tablero general & Consultar un resumen del estado del sistema al
		entrar a la consola. \\ \hline
	
		RF-A03 & Gestionar cuentas de personal & Listar, crear con contraseña temporal, editar,
		restablecer contraseña, dar de baja y reactivar cuentas de cualquiera de los cinco roles.
		\\ \hline
	
		RF-A04 & Administrar el catálogo institucional & Listar, crear y editar dependencias, e
		importar o actualizar el catálogo de forma masiva desde un archivo de hoja de cálculo.
		\\ \hline
	
		RF-A05 & Administrar plantillas de documentos & Listar plantillas, consultar los marcadores
		disponibles, previsualizar con datos de ejemplo y editar o publicar su contenido. \\ \hline
	
		RF-A06 & Gestionar respaldos de la base de datos & Listar los respaldos existentes, ejecutar
		un respaldo bajo demanda, descargar un respaldo y restaurar la base de datos a partir de uno
		de ellos. \\ \hline
	
		RF-A07 & Consultar la bitácora de acciones críticas & Consultar el registro de acciones
		administrativas con usuario, acción, dirección de origen y fecha. \\ \hline
	
		RF-A08 & Administrar el contenido de preguntas frecuentes & Listar, crear, editar y eliminar
		las preguntas del portal público. \\ \hline
	
		RF-A09 & Cambiar la propia contraseña & Modificar la contraseña de la cuenta propia desde la
		consola. \\ \hline
	\end{longtable}
	
	\normalsize
	\setlength{\tabcolsep}{6pt}
	\renewcommand{\arraystretch}{1}
	
	\subsubsection{Requerimientos no funcionales}
	Los requerimientos no funcionales expresan atributos de calidad y no funciones. Se agruparon
	por la categoría del atributo que abordan ---disponibilidad, confidencialidad, resiliencia,
	usabilidad y trazabilidad--- porque esa clasificación es la que permite verificarlos: cada
	categoría se comprueba con un tipo distinto de prueba. La Tabla~\ref{tab:rnf} los concentra.
	
	\rowcolors{2}{rowblue}{white}
	
	\small
	\renewcommand{\arraystretch}{1.2}
	\setlength{\tabcolsep}{4pt}
	
	\begin{longtable}{|p{1.7cm}|p{2.7cm}|p{9.7cm}|}
		\caption{Requerimientos no funcionales de la plataforma}
		\label{tab:rnf} \\
		\hline
		\rowcolor{headerblue}
		\headercell{ID} & \headercell{Categoría} & \headercell{Atributo requerido} \\
		\hline
		\endfirsthead
	
		\hline
		\rowcolor{headerblue}
		\headercell{ID} & \headercell{Categoría} & \headercell{Atributo requerido} \\
		\hline
		\endhead
	
		RNF-Q01 & Disponibilidad & Debe poder presentarse una queja sin cuenta ni sesión, en
		cualquier momento y desde cualquier dispositivo con navegador. \\ \hline
	
		RNF-Q02 & Confidencialidad & Todo el tráfico público debe viajar cifrado con TLS; tras el
		inicio de sesión, la autenticación se transmite como token en el encabezado de autorización y
		nunca como credenciales en claro. \\ \hline
	
		RNF-Q03 & Capacidad & El proxy público debe aceptar peticiones de hasta 100~MB y el backend
		hasta 60~MB por petición multiparte, para permitir varios archivos de evidencia por queja.
		\\ \hline
	
		RNF-Q04 & Resiliencia & Un fallo en el envío de correo o en el registro de una notificación
		nunca debe impedir que la operación principal se complete; el fallo se degrada a un registro
		en la bitácora del servidor. \\ \hline
	
		RNF-Q05 & Usabilidad & Los formularios extensos deben presentarse en pasos numerados, con
		previsualización de los archivos seleccionados y ayuda contextual que no reduzca el espacio
		útil del formulario. \\ \hline
	
		RNF-Q06 & Seguridad & La sesión del quejoso debe expirar en una hora; al expirar o ser
		rechazada, el cliente debe cerrar sesión de forma automática. \\ \hline
	
		RNF-Q07 & Usabilidad & Toda operación debe producir retroalimentación visual inmediata de
		éxito, error o validación; no se admiten fallos silenciosos. \\ \hline
	
		RNF-R01 & Trazabilidad & Cada rechazo y cada turnado debe registrar, sin excepción, el
		correo de quien lo procesó y la fecha en que ocurrió. \\ \hline
	
		RNF-R02 & Seguridad & El personal debe autenticarse en un único punto; ningún panel de
		personal implementa lógica de contraseñas propia. \\ \hline
	
		RNF-R03 & Interoperabilidad & La exportación del historial debe producir un archivo de hoja
		de cálculo en formato estándar, abrible sin conversión previa. \\ \hline
	
		RNF-R04 & Resiliencia & Si falla el envío del correo de rechazo, el rechazo en sí debe
		completarse igualmente. \\ \hline
	
		RNF-R05 & Consistencia & Los paneles de personal deben compartir componentes visuales,
		patrones de error y criterios de accesibilidad. \\ \hline
	
		RNF-A01 & Seguridad & El primer acceso administrativo debe autoprovisionarse en el arranque
		inicial con una contraseña temporal registrada una sola vez en la bitácora del servidor;
		nunca incrustada en el cliente ni expuesta por la interfaz de programación. \\ \hline
	
		RNF-A02 & Trazabilidad & Toda acción administrativa relevante, no solo el inicio de sesión,
		debe quedar registrada con la dirección de origen. \\ \hline
	
		RNF-A03 & Seguridad & La restauración de la base de datos debe exigir una confirmación
		explícita en el cuerpo de la petición; no puede dispararse con una sola acción accidental.
		\\ \hline
	
		RNF-A04 & Tolerancia a fallos & La importación masiva del catálogo debe reportar un resumen
		de filas procesadas, creadas, actualizadas y con error, en lugar de fallar en bloque ante una
		fila mal formada. \\ \hline
	
		RNF-A05 & Usabilidad & La consola de administración debe aprovechar la densidad de
		información disponible en pantalla. \\ \hline
	
		RNF-A06 & Seguridad & La generación, descarga y restauración de respaldos debe estar
		disponible únicamente para el rol de administración de sistemas. \\ \hline
	\end{longtable}
	
	\normalsize
	
	\subsubsection{Reglas de negocio}
	Las reglas de negocio son las restricciones que el sistema debe hacer cumplir para que el
	procedimiento sea válido. A diferencia de los requerimientos funcionales, que describen lo que
	el sistema \emph{permite} hacer, las reglas describen lo que \emph{prohíbe} o \emph{obliga}. Su
	característica común en esta plataforma es que se validan en el servidor y no solo en la
	interfaz: ocultar un botón no es hacer cumplir una regla, y así se hace explícito en varias de
	las que siguen.
	
	\small
	\renewcommand{\arraystretch}{1.2}
	
	\begin{longtable}{|p{1.7cm}|p{2.7cm}|p{9.7cm}|}
		\caption{Reglas de negocio de la plataforma}
		\label{tab:reglas_negocio} \\
		\hline
		\rowcolor{headerblue}
		\headercell{ID} & \headercell{Temática} & \headercell{Regla} \\
		\hline
		\endfirsthead
	
		\hline
		\rowcolor{headerblue}
		\headercell{ID} & \headercell{Temática} & \headercell{Regla} \\
		\hline
		\endhead
	
		RN-Q01 & Menores de edad & Una queja de una persona menor de 18 años requiere datos de tutor
		o adulto responsable capturados y confirmados antes de poder enviarse. \\ \hline
	
		RN-Q02 & Menores de edad & Una persona menor de 14 años no puede autorregistrarse ni
		gestionar la queja por sí misma; la interfaz indica de forma explícita que debe ser el tutor
		quien complete el formulario a su nombre. \\ \hline
	
		RN-Q03 & Identificación & El número de boleta o de empleado debe ser exclusivamente numérico
		y no exceder doce caracteres. \\ \hline
	
		RN-Q04 & Activación & Una cuenta de quejoso solo puede activarse si existe al menos una queja
		previa con ese folio \emph{y} ese correo; el par folio-correo es la llave de activación.
		\\ \hline
	
		RN-Q05 & Activación & Una cuenta ya activa no puede volver a activarse. \\ \hline
	
		RN-Q06 & Contraseñas & La contraseña debe tener al menos ocho caracteres, una mayúscula y un
		dígito, tanto en la activación de cuenta como en el restablecimiento. \\ \hline
	
		RN-Q07 & Contraseñas & Un código de recuperación expira diez minutos después de solicitado y
		solo es válido una vez. \\ \hline
	
		RN-Q08 & Edición & Una queja solo es editable por su autor mientras su estatus sea
		\texttt{RECIBIDA}; en cualquier otro estatus la edición se rechaza en el servidor, no
		únicamente se oculta en la interfaz. \\ \hline
	
		RN-Q09 & Conciliación & Un acuerdo de conciliación solo puede ser respondido por el quejoso
		al que va dirigido; el correo del acuerdo debe coincidir con el del token de sesión.
		\\ \hline
	
		RN-Q10 & Origen & El origen de una queja se asigna de forma automática según el punto de
		acceso empleado ---\texttt{PUBLICO}, \texttt{AUTENTICADO} o \texttt{MANUAL}---; no es un dato
		que el usuario declare. \\ \hline
	
		RN-R01 & Ciclo de vida & La queja transita \texttt{RECIBIDA} $\rightarrow$
		\texttt{EN\_VALIDACION} $\rightarrow$ \texttt{RECHAZADA} o \texttt{TURNADA}. Una vez
		rechazada o turnada no regresa a un estado anterior desde el panel de revisión. \\ \hline
	
		RN-R02 & Autorización & Únicamente las cuentas con rol de recepción pueden invocar los
		puntos de acceso del servicio de revisión; la restricción se declara en cada controlador y no
		solo a nivel de ruta. \\ \hline
	
		RN-R03 & Rechazo & Rechazar una queja exige al menos un motivo y genera de forma automática
		el correo de rechazo al quejoso; no es una acción silenciosa. \\ \hline
	
		RN-R04 & Turnado & Turnar una queja exige área y defensor asignado, y genera un folio oficial
		de turnado distinto del folio original de seguimiento. \\ \hline
	
		RN-R05 & Registro manual & Una queja registrada a partir de un documento físico recibe origen
		\texttt{MANUAL} y admite que no se conozca de antemano la boleta o el número de empleado.
		\\ \hline
	
		RN-R06 & Turnado & El catálogo de personal disponible para turnar incluye únicamente cuentas
		activas con rol de defensor o subdefensor. \\ \hline
	
		RN-R07 & Consistencia & El servicio de quejas y el de revisión operan sobre la misma tabla
		física mediante entidades independientes, de modo que un cambio de estatus hecho por
		recepción es visible de inmediato para el quejoso sin necesidad de una llamada entre
		servicios. \\ \hline
	
		RN-A01 & Autorización & Todos los puntos de acceso administrativos exigen el rol de
		administración de sistemas, salvo los de autenticación y los de gestión del perfil propio,
		que cualquier cuenta autenticada puede usar sobre sí misma. \\ \hline
	
		RN-A02 & Cuentas & Una cuenta nueva de personal se crea siempre como temporal y con
		obligación de cambiar la contraseña en el primer inicio de sesión. \\ \hline
	
		RN-A03 & Cuentas & Dar de baja una cuenta la desactiva, no la elimina, para conservar la
		integridad referencial con las acciones que registró en el pasado. \\ \hline
	
		RN-A04 & Unicidad & El correo institucional y el número de empleado del personal son únicos
		en el sistema, garantizado por restricción de base de datos y no solo por validación de
		interfaz. \\ \hline
	
		RN-A05 & Autorización & Las rutas de administración de un catálogo deben evaluarse antes que
		la regla pública del mismo prefijo; sin ese orden explícito prevalecería la regla menos
		restrictiva por ser más genérica. \\ \hline
	
		RN-A06 & Plantillas & Una plantilla de documento se identifica por una clave estable de
		negocio y no por su identificador autogenerado, para poder referenciarla sin depender del
		orden de creación. \\ \hline
	
		RN-A07 & Trazabilidad & Cada actualización del contenido de una plantilla registra quién la
		actualizó y cuándo, además del asiento correspondiente en la bitácora. \\ \hline
	\end{longtable}
	
	\normalsize
	\setlength{\tabcolsep}{6pt}
	\renewcommand{\arraystretch}{1}
	
	\subsection{Casos de uso}
	\label{sec:casos_uso}
	Los casos de uso describen las interacciones posibles entre cada actor y el sistema con
	independencia de la tecnología empleada. Se presentan agrupados por actor y con la estructura
	mínima que permite verificarlos: precondición, flujo principal y postcondición. La especificación
	extendida de cada caso, con flujos alternos y de excepción, se documenta en el anexo
	correspondiente.
	
	\small
	\renewcommand{\arraystretch}{1.2}
	\setlength{\tabcolsep}{4pt}
	
	\begin{longtable}{|p{1.5cm}|p{2.8cm}|p{2.6cm}|p{3.4cm}|p{3.4cm}|}
		\caption{Casos de uso del actor Quejoso}
		\label{tab:cu_quejoso} \\
		\hline
		\rowcolor{headerblue}
		\headercell{ID} & \headercell{Nombre} & \headercell{Precondición} &
		\headercell{Flujo principal} & \headercell{Postcondición} \\
		\hline
		\endfirsthead
	
		\hline
		\rowcolor{headerblue}
		\headercell{ID} & \headercell{Nombre} & \headercell{Precondición} &
		\headercell{Flujo principal} & \headercell{Postcondición} \\
		\hline
		\endhead
	
		CU-Q01 & Presentar una queja sin cuenta & Ninguna & Completa el formulario público con
		identidad, motivo, evidencias y tutor si aplica, y lo envía & Queja creada con estatus
		\texttt{RECIBIDA} y origen \texttt{PUBLICO}; se muestra el folio \\ \hline
	
		CU-Q02 & Consultar una queja por folio & Conoce folio y correo & Captura ambos datos en la
		pantalla de consulta & Se muestra el detalle y el estatus de la queja \\ \hline
	
		CU-Q03 & Activar cuenta de seguimiento & Existe una queja con ese folio y correo & Define una
		contraseña; el sistema valida el par folio-correo contra el servicio de quejas & Cuenta
		activa; puede iniciar sesión \\ \hline
	
		CU-Q04 & Iniciar sesión & Cuenta activa & Captura correo y contraseña & Token de sesión
		emitido; se registra la notificación de acceso \\ \hline
	
		CU-Q05 & Recuperar contraseña & Cuenta activa existente & Solicita el código, lo recibe por
		correo y lo captura con la nueva contraseña & Contraseña actualizada; se envía correo de
		confirmación \\ \hline
	
		CU-Q06 & Presentar una queja autenticado & Sesión iniciada & Desde el panel captura motivo,
		descripción y evidencias & Queja creada con estatus \texttt{RECIBIDA} y origen
		\texttt{AUTENTICADO} \\ \hline
	
		CU-Q07 & Editar una queja propia & Sesión iniciada y queja en \texttt{RECIBIDA} & Modifica
		los campos habilitados desde el detalle y guarda & Queja actualizada; en otro estatus el
		servidor rechaza la edición \\ \hline
	
		CU-Q08 & Responder un acuerdo de conciliación & Existe un acuerdo pendiente dirigido a su
		correo & Revisa el acuerdo y lo acepta o rechaza con comentario opcional & Acuerdo pasa a
		aceptado o rechazado, con fecha de respuesta \\ \hline
	
		CU-Q09 & Revisar notificaciones & Sesión iniciada & Abre el centro de notificaciones y marca
		como leídas & El contador de no leídas se actualiza \\ \hline
	
		CU-Q10 & Actualizar el perfil & Sesión iniciada & Edita los campos habilitados y guarda &
		Perfil actualizado \\ \hline
	\end{longtable}
	
	\begin{longtable}{|p{1.5cm}|p{2.8cm}|p{2.6cm}|p{3.4cm}|p{3.4cm}|}
		\caption{Casos de uso del actor Personal de recepción}
		\label{tab:cu_recepcion} \\
		\hline
		\rowcolor{headerblue}
		\headercell{ID} & \headercell{Nombre} & \headercell{Precondición} &
		\headercell{Flujo principal} & \headercell{Postcondición} \\
		\hline
		\endfirsthead
	
		\hline
		\rowcolor{headerblue}
		\headercell{ID} & \headercell{Nombre} & \headercell{Precondición} &
		\headercell{Flujo principal} & \headercell{Postcondición} \\
		\hline
		\endhead
	
		CU-R01 & Iniciar sesión en el panel de revisión & Cuenta activa con rol de recepción &
		Se autentica contra el servicio de administración & Token emitido; acceso al panel \\ \hline
	
		CU-R02 & Revisar la bandeja de entrada & Sesión iniciada & Consulta contadores y lista de
		pendientes & --- \\ \hline
	
		CU-R03 & Validar y rechazar una queja & Queja en \texttt{RECIBIDA} o \texttt{EN\_VALIDACION} &
		Revisa documentos y antecedentes, marca motivos y observaciones, confirma & Queja pasa a
		\texttt{RECHAZADA}; se envía correo al quejoso \\ \hline
	
		CU-R04 & Validar y turnar una queja & Queja en \texttt{RECIBIDA} o \texttt{EN\_VALIDACION} &
		Revisa antecedentes, elige área y defensor, agrega comentarios, confirma & Queja pasa a
		\texttt{TURNADA} con folio oficial, área y defensor \\ \hline
	
		CU-R05 & Registrar una queja en papel & Documento físico recibido & Captura los datos del
		documento y adjunta el escaneo si lo hay & Queja creada con origen \texttt{MANUAL} y folio
		propio \\ \hline
	
		CU-R06 & Consultar y exportar el historial & Existen trámites procesados & Filtra por texto,
		estatus o fecha y exporta si lo requiere & Archivo de hoja de cálculo descargado con el
		filtro aplicado \\ \hline
	
		CU-R07 & Emitir un acuerdo de conciliación & Existe una queja con folio conocido & Captura
		asunto y términos dirigidos al folio y correo del quejoso & Acuerdo creado como pendiente,
		visible para el quejoso \\ \hline
	\end{longtable}
	
	\begin{longtable}{|p{1.5cm}|p{2.8cm}|p{2.6cm}|p{3.4cm}|p{3.4cm}|}
		\caption{Casos de uso del actor Administrador de TI}
		\label{tab:cu_admin} \\
		\hline
		\rowcolor{headerblue}
		\headercell{ID} & \headercell{Nombre} & \headercell{Precondición} &
		\headercell{Flujo principal} & \headercell{Postcondición} \\
		\hline
		\endfirsthead
	
		\hline
		\rowcolor{headerblue}
		\headercell{ID} & \headercell{Nombre} & \headercell{Precondición} &
		\headercell{Flujo principal} & \headercell{Postcondición} \\
		\hline
		\endhead
	
		CU-A01 & Iniciar sesión como administrador & Cuenta activa de administración &
		Se autentica en la consola & Token emitido; acceso al panel \\ \hline
	
		CU-A02 & Dar de alta personal & Sesión iniciada & Captura nombre, número de empleado, correo
		y rol & Cuenta creada con contraseña temporal; queda en bitácora \\ \hline
	
		CU-A03 & Restablecer la contraseña de otra cuenta & La cuenta destino existe & Selecciona la
		cuenta y solicita el restablecimiento & Nueva contraseña temporal generada; queda en bitácora
		\\ \hline
	
		CU-A04 & Dar de baja o reactivar una cuenta & Sesión iniciada & Selecciona la cuenta y
		confirma la operación & La cuenta cambia su estado de actividad; queda en bitácora \\ \hline
	
		CU-A05 & Importar el catálogo institucional & Archivo de hoja de cálculo con el formato
		esperado & Carga el archivo en el módulo de catálogo & Dependencias creadas o actualizadas
		de forma masiva; se muestra el resumen \\ \hline
	
		CU-A06 & Editar una plantilla de oficio & Sesión iniciada & Edita el contenido con los
		marcadores disponibles, previsualiza y publica & Plantilla actualizada; queda en bitácora
		\\ \hline
	
		CU-A07 & Ejecutar un respaldo & Sesión iniciada & Solicita un respaldo bajo demanda &
		Archivo de respaldo generado y listado; queda en bitácora \\ \hline
	
		CU-A08 & Restaurar desde un respaldo & Existe al menos un respaldo & Selecciona el archivo y
		confirma de forma explícita la operación destructiva & Base de datos restaurada; queda en
		bitácora \\ \hline
	
		CU-A09 & Revisar la bitácora de seguridad & Sesión iniciada & Abre el historial de acciones
		críticas & Se listan las acciones con usuario, dirección de origen y fecha \\ \hline
	
		CU-A10 & Administrar las preguntas frecuentes & Sesión iniciada & Crea, edita o elimina
		preguntas & Los cambios se reflejan de inmediato en el portal público \\ \hline
	\end{longtable}
	
	\normalsize
	\setlength{\tabcolsep}{6pt}
	\renewcommand{\arraystretch}{1}
	
	\subsection{Ciclo de vida del expediente}
	\label{sec:maquina_estados}
	El elemento que ordena todo el flujo de trabajo es el estatus de la queja. Se modeló como una
	máquina de estados con cuatro valores y transiciones dirigidas, decisión que se tomó por dos
	razones: permite validar en el servidor qué operaciones son admisibles en cada momento ---la
	regla RN-Q08 depende por completo de esto--- y da al quejoso una línea de tiempo comprensible
	del avance de su expediente.
	
	\begin{itemize}
		\item \texttt{RECIBIDA}. Estado inicial de toda queja, sin importar su origen. Es el único
		estado en el que el quejoso puede editar su propia queja, porque es el único en el que la
		Defensoría no ha empezado a trabajar sobre ella.
	
		\item \texttt{EN\_VALIDACION}. El personal de recepción tomó la queja y está verificando su
		documentación y sus antecedentes. La edición por parte del quejoso queda cerrada a partir de
		aquí, para evitar que el expediente cambie mientras se lo evalúa.
	
		\item \texttt{RECHAZADA}. Estado terminal. La queja no reunió los requisitos de
		documentación; el rechazo exige motivos registrados y notificación al quejoso.
	
		\item \texttt{TURNADA}. Estado terminal desde la perspectiva de recepción. La queja fue
		canalizada a un área y a un defensor, y recibió un folio oficial de turnado.
	\end{itemize}
	
	Los dos estados terminales lo son únicamente respecto del alcance implementado en esta etapa:
	el flujo completo de la Defensoría continúa después del turnado con la investigación, la
	conciliación y la resolución, etapas que corresponden a los roles que aún no cuentan con
	interfaz propia. El modelo de datos ya contempla los campos necesarios para esas etapas, de modo
	que su incorporación no requerirá modificar el esquema existente.
	
	\subsection{Diseño de la arquitectura}
	\label{sec:diseno_arquitectura}
	La arquitectura de la plataforma aplica los principios expuestos en la
	Sección~\ref{sec:arquitectura_software}. Esta sección describe cómo se concretaron: qué
	componentes existen, cómo se comunican, cómo se despliegan y qué alternativas se descartaron en
	cada decisión.
	
	\subsubsection{Vista de componentes}
	El sistema se descompuso en siete microservicios, cada uno con una responsabilidad acotada y su
	propio ciclo de despliegue. El criterio de división no fue técnico sino de dominio: cada
	servicio agrupa las operaciones que comparten un mismo conjunto de datos y una misma razón para
	cambiar. La Tabla~\ref{tab:microservicios} presenta la descomposición resultante.
	
	\rowcolors{2}{rowblue}{white}
	
	\small
	\renewcommand{\arraystretch}{1.25}
	\setlength{\tabcolsep}{4pt}
	
	\begin{longtable}{|p{3.4cm}|p{1.2cm}|p{6.6cm}|p{2.4cm}|}
		\caption{Microservicios de la plataforma y su responsabilidad}
		\label{tab:microservicios} \\
		\hline
		\rowcolor{headerblue}
		\headercell{Servicio} & \headercell{Puerto} & \headercell{Responsabilidad} &
		\headercell{Emite tokens} \\
		\hline
		\endfirsthead
	
		\hline
		\rowcolor{headerblue}
		\headercell{Servicio} & \headercell{Puerto} & \headercell{Responsabilidad} &
		\headercell{Emite tokens} \\
		\hline
		\endhead
	
		\texttt{auth-service} & 8083 & Sesión del quejoso, activación de cuenta, recuperación de
		contraseña y perfil & Sí, para quejosos \\ \hline
	
		\texttt{queja-service} & 8084 & Registro de quejas público y autenticado, consulta y edición
		de quejas propias, evidencias y respuesta a acuerdos de conciliación & No, solo verifica
		\\ \hline
	
		\texttt{notificaciones-\allowbreak service} & 8085 & Envío de correo electrónico y centro de
		notificaciones persistido por usuario & No, solo verifica \\ \hline
	
		\texttt{catalogo-\allowbreak service} & 8086 & Catálogo de dependencias y unidades académicas del
		Instituto & No, solo verifica \\ \hline
	
		\texttt{admin-service} & 8087 & Sesión del personal, cuentas y roles, plantillas de oficios,
		respaldos y bitácora de seguridad & Sí, para personal \\ \hline
	
		\texttt{revision-\allowbreak service} & 8088 & Bandeja de recepción, validación, rechazo, turnado,
		registro manual, historial exportable y emisión de acuerdos & No, solo verifica \\ \hline
	
		\texttt{chatbot-\allowbreak service} & 8089 & Menú de preguntas frecuentes del portal público y su
		administración & No, solo verifica \\ \hline
	\end{longtable}
	
	\normalsize
	
	Solo dos de los siete servicios emiten tokens de sesión: uno para quejosos y otro para
	personal. Los cinco restantes se limitan a verificarlos. Esta asimetría es deliberada y
	responde al principio de un único punto de emisión de credenciales: concentrar la lógica de
	contraseñas en dos lugares en vez de siete reduce la superficie donde puede introducirse un
	defecto de autenticación, y evita que la política de complejidad de contraseñas quede
	implementada de siete maneras distintas.
	
	\subsubsection{Vista de despliegue}
	El sistema se despliega sobre dos servidores virtuales privados con responsabilidades separadas.
	El primero, al que en adelante se denomina \emph{servidor de aplicación}, aloja los siete
	microservicios y la base de datos. El segundo, el \emph{servidor de acceso público}, aloja las
	tres aplicaciones web y el proxy inverso que da la cara a Internet.
	
	Esta separación cumple tres propósitos. El primero es de seguridad: el servidor que contiene la
	base de datos y la lógica de negocio no está expuesto directamente a Internet, y su cortafuegos
	acepta conexiones a los puertos de los microservicios únicamente desde la dirección del servidor
	de acceso público. El segundo es de aislamiento de fallos: una saturación en la entrega de
	archivos estáticos no compite por recursos con el procesamiento de las peticiones de negocio.
	El tercero es operativo: el certificado de transporte y la configuración del dominio residen en
	un solo lugar, de modo que renovarlos no implica reiniciar ningún servicio de negocio.
	
	El único puerto realmente público es el de HTTPS; el puerto de HTTP se conserva abierto
	exclusivamente para redirigir el tráfico a su equivalente cifrado y para atender el reto de
	validación del certificado. La Tabla~\ref{tab:enrutamiento} presenta las reglas de enrutamiento
	del proxy inverso, que son el contrato entre el exterior y la arquitectura interna.
	
	\small
	\renewcommand{\arraystretch}{1.2}
	
	\begin{longtable}{|p{4.0cm}|p{5.4cm}|p{4.4cm}|}
		\caption{Reglas de enrutamiento del proxy inverso público}
		\label{tab:enrutamiento} \\
		\hline
		\rowcolor{headerblue}
		\headercell{Ruta pública} & \headercell{Destino} & \headercell{Servidor} \\
		\hline
		\endfirsthead
	
		\hline
		\rowcolor{headerblue}
		\headercell{Ruta pública} & \headercell{Destino} & \headercell{Servidor} \\
		\hline
		\endhead
	
		\texttt{/} & Portal del quejoso & Acceso público \\ \hline
		\texttt{/admin/} & Consola de administración & Acceso público \\ \hline
		\texttt{/revision/} & Panel de revisión & Acceso público \\ \hline
		\texttt{/api/auth/} & \texttt{auth-service} & Aplicación \\ \hline
		\texttt{/api/quejoso/} & \texttt{queja-service} & Aplicación \\ \hline
		\texttt{/api/notificaciones/} & \texttt{notificaciones-\allowbreak service} & Aplicación \\ \hline
		\texttt{/api/catalogos/} & \texttt{catalogo-\allowbreak service} & Aplicación \\ \hline
		\texttt{/api/admin/} & \texttt{admin-service} & Aplicación \\ \hline
		\texttt{/api/revision/} & \texttt{revision-\allowbreak service} & Aplicación \\ \hline
		\texttt{/api/chatbot/} & \texttt{chatbot-\allowbreak service} & Aplicación \\ \hline
	\end{longtable}
	
	\normalsize
	\setlength{\tabcolsep}{6pt}
	\renewcommand{\arraystretch}{1}
	
	Una consecuencia de diseño de este esquema merece destacarse: como el proxy entrega tanto las
	aplicaciones web como la interfaz de programación bajo un mismo dominio, el navegador considera
	que todo proviene del mismo origen y no activa la verificación de intercambio de recursos entre
	orígenes. Eso simplifica el modelo de seguridad del cliente sin sacrificar la separación en
	servicios del lado del servidor.
	
	\subsubsection{Diseño de la comunicación entre servicios}
	La comunicación entre servicios es síncrona sobre HTTP siguiendo el estilo REST, conforme a lo
	expuesto en la Sección~\ref{sec:arquitectura_software}. Se descartó la mensajería asíncrona
	porque el volumen de operaciones y el número de servicios no justifican la complejidad operativa
	de un intermediario de mensajes, y porque ninguno de los flujos identificados tolera bien la
	consistencia eventual: un quejoso que envía una queja espera su folio en la misma respuesta, no
	después.
	
	El diseño reduce al mínimo las llamadas entre servicios. Solo existen dos rutas de comunicación
	interna: la validación del par folio-correo que el servicio de autenticación solicita al de
	quejas durante la activación de una cuenta, y el envío de correo que varios servicios solicitan
	al de notificaciones. Todo lo demás se resuelve por lectura directa de la base de datos
	compartida, decisión cuyas consecuencias se discuten en la
	Sección~\ref{sec:decisiones_diseno}.
	
	\subsubsection{Diseño de la seguridad y el control de acceso}
	El esquema de seguridad se construyó sobre cuatro mecanismos, cada uno atendiendo una amenaza
	distinta.
	
	La \textbf{autenticación} emplea tokens firmados con algoritmo simétrico. Un token emitido por
	el servicio de autenticación es válido en los servicios que atienden al quejoso, y uno emitido
	por el servicio de administración es válido en los que atienden al personal, porque todos
	comparten la misma clave de firma. Cada servicio verifica el token de forma local, sin consultar
	al emisor, lo que evita un punto único de falla en cada petición. La vigencia del token es de
	una hora, y el cliente cierra la sesión de forma automática cuando el servidor la rechaza.
	
	El \textbf{almacenamiento de contraseñas} usa una función de resumen con costo configurable en
	los dos servicios que emiten sesión. La política de complejidad exige un mínimo de ocho
	caracteres, al menos una mayúscula y al menos un dígito, y está implementada como una expresión
	compartida para que no divergja entre servicios. Las cuentas de personal se crean con contraseña
	temporal y obligación de cambiarla en el primer acceso.
	
	La \textbf{autorización} se declara por rol en cada controlador, no solo por ruta. Esta
	distinción no es cosmética: una restricción declarada únicamente en la configuración de rutas
	puede eludirse si una ruta nueva se agrega sin actualizar esa configuración, mientras que la
	declaración en el controlador viaja con el código que protege.
	
	El \textbf{cifrado en tránsito} se termina en el proxy inverso con un certificado de una
	autoridad reconocida y renovación automática. Del proxy hacia los microservicios el tráfico
	viaja sin cifrar, decisión aceptable porque ese tramo no atraviesa Internet y está restringido
	por dirección de origen en el cortafuegos, pero que conviene revisar si en el futuro los
	servidores dejan de estar bajo control directo del proyecto.
	
	Un punto pendiente que debe declararse: el cierre automático de sesión por inactividad
	solicitado para las cuentas de personal no está implementado al cierre de esta etapa. La
	expiración del token limita la ventana de exposición a una hora, pero no equivale a un cierre
	por inactividad.
	
	\subsubsection{Decisiones de diseño y alternativas descartadas}
	\label{sec:decisiones_diseno}
	Esta sección documenta las decisiones estructurales del proyecto con el formato
	contexto-alternativas-decisión-consecuencias. Registrarlas de esta forma cumple un propósito
	concreto: una decisión de arquitectura cuyo razonamiento no se documenta se vuelve
	indistinguible de un accidente, y quien mantenga el sistema en el futuro no sabrá si puede
	revertirla.
	
	\paragraph{Decisión 1: base de datos única compartida.}
	\emph{Contexto:} siete servicios independientes necesitan datos que se relacionan entre sí ---una
	queja, su tutor, sus evidencias, el usuario que la presentó, la notificación que se le envió---.
	\emph{Alternativas:} una base de datos por servicio, con sincronización explícita entre ellas;
	o una base única compartida. \emph{Decisión:} base única. \emph{Consecuencias aceptadas:} se
	sacrifica la autonomía de datos que caracteriza a los microservicios puros, y dos servicios
	escriben la misma tabla sin coordinación transaccional entre ellos. A cambio se evita la
	complejidad de mantener consistencia entre bases separadas, que a la escala de este sistema
	habría requerido un esfuerzo mayor que el beneficio obtenido. Es un híbrido consciente, no una
	implementación incompleta del patrón.
	
	\paragraph{Decisión 2: entidades independientes sobre una tabla compartida.}
	\emph{Contexto:} el servicio de quejas y el de revisión necesitan la misma tabla de quejas, pero
	con vistas distintas: el primero requiere los campos del quejoso, el segundo los del flujo de
	validación. \emph{Alternativas:} un servicio dueño de la tabla que exponga los datos al otro por
	interfaz de programación; o dos clases de persistencia independientes sobre la misma tabla
	física. \emph{Decisión:} entidades independientes. \emph{Consecuencias aceptadas:} un cambio de
	estatus hecho por recepción es visible de inmediato para el quejoso sin llamada entre servicios,
	lo que simplifica el flujo; pero ambos servicios pueden agregar columnas sin que el otro lo
	sepa, y una modificación descoordinada del esquema puede romper al servicio vecino. La mitigación
	adoptada es que cada servicio declare únicamente las columnas que usa.
	
	\paragraph{Decisión 3: catálogo institucional como servicio propio.}
	\emph{Contexto:} el catálogo de dependencias del Instituto es requerido por el formulario público
	de quejas, por el turnado en el panel de revisión y por la consola de administración.
	\emph{Alternativas:} alojarlo dentro del servicio de quejas, que fue quien primero lo necesitó;
	o extraerlo a un servicio propio. \emph{Decisión:} servicio propio.
	\emph{Consecuencias aceptadas:} se añadió un servicio más que operar y desplegar, con el
	sobrecosto correspondiente, en un momento en que un solo consumidor lo necesitaba. La razón para
	asumirlo fue anticipada: al convertirse en dato de referencia de tres consumidores distintos,
	alojarlo dentro de un servicio de negocio habría obligado a que los otros dos dependieran de él
	por razones ajenas a su dominio.
	
	\paragraph{Decisión 4: evidencias almacenadas en la base de datos.}
	\emph{Contexto:} cada queja puede llevar varios archivos de evidencia que constituyen material
	probatorio. \emph{Alternativas:} guardar los archivos en el sistema de archivos del servidor y
	almacenar la ruta en la base; o almacenar el contenido binario completo dentro de la base.
	\emph{Decisión:} contenido binario en la base. \emph{Consecuencias aceptadas:} la base crece más
	rápido y los respaldos son más grandes. A cambio se obtiene lo que en este dominio pesa más: el
	archivo y su expediente se respaldan y se restauran como una unidad atómica, de modo que no puede
	existir un expediente cuya evidencia se haya perdido porque el respaldo del sistema de archivos y
	el de la base se tomaron en momentos distintos. Para material con posible valor probatorio, esa
	garantía justifica el costo de almacenamiento.
	
	\paragraph{Decisión 5: tres aplicaciones web independientes.}
	\emph{Contexto:} tres actores con perfiles de uso y niveles de privilegio muy distintos.
	\emph{Alternativas:} una sola aplicación con vistas y rutas condicionadas por rol; o tres
	aplicaciones separadas. \emph{Decisión:} tres aplicaciones. \emph{Consecuencias aceptadas:} se
	duplica cierto código de infraestructura ---guardas, interceptores, estilos base--- en tres
	proyectos que deben mantenerse sincronizados a mano. A cambio, el código de la consola de
	administración no se entrega al navegador de un usuario público, lo que reduce la superficie
	expuesta: un visitante del portal no recibe siquiera las rutas de administración.
	
	\paragraph{Decisión 6: datos del quejoso en columnas propias.}
	\emph{Contexto:} el registro público debe capturar la identidad completa del quejoso, que en una
	versión inicial se concatenaba como texto libre dentro de la descripción de la queja.
	\emph{Alternativas:} conservar la concatenación, más simple de implementar; o estructurar cada
	dato en su propia columna. \emph{Decisión:} columnas propias. \emph{Consecuencias aceptadas:} el
	esquema creció y el formulario requiere más validación. A cambio, los datos se volvieron
	consultables y verificables ---la búsqueda de antecedentes del requerimiento RF-R04 sería
	impracticable sobre texto concatenado--- y la unidad académica pudo pasar de texto libre a una
	referencia al catálogo institucional.
	
	\paragraph{Decisión 7: una imagen de contenedor por servicio.}
	\emph{Contexto:} los microservicios comparten la misma imagen base y el mismo procedimiento de
	construcción. \emph{Alternativas:} una imagen común reutilizada por todos, parametrizada al
	ejecutarse; o una imagen dedicada por servicio. \emph{Decisión:} una por servicio.
	\emph{Consecuencias aceptadas:} se consume más espacio de almacenamiento en el registro de
	imágenes. A cambio, cada contenedor es identificable y versionable de forma independiente, y se
	elimina la ambigüedad de no poder saber qué versión de qué servicio está ejecutando un
	contenedor determinado.
	
	\subsection{Diseño de la base de datos}
	\label{sec:diseno_bd}
	La base de datos es relacional y única para los siete microservicios. Cada tabla tiene un
	servicio propietario, entendido como el único autorizado a definir su estructura, y en dos casos
	una segunda tabla es escrita por un servicio adicional bajo el patrón descrito en la Decisión~2.
	La Tabla~\ref{tab:esquema_bd} presenta el esquema con su propietario y su contenido.
	
	\rowcolors{2}{rowblue}{white}
	
	\small
	\renewcommand{\arraystretch}{1.25}
	\setlength{\tabcolsep}{4pt}
	
	\begin{longtable}{|p{3.5cm}|p{3.1cm}|p{7.5cm}|}
		\caption{Esquema de la base de datos por tabla y servicio propietario}
		\label{tab:esquema_bd} \\
		\hline
		\rowcolor{headerblue}
		\headercell{Tabla} & \headercell{Propietario} & \headercell{Contenido} \\
		\hline
		\endfirsthead
	
		\hline
		\rowcolor{headerblue}
		\headercell{Tabla} & \headercell{Propietario} & \headercell{Contenido} \\
		\hline
		\endhead
	
		\texttt{usuarios} & \texttt{auth-service} & Cuentas de quejoso: identidad institucional,
		credencial cifrada, datos de contacto, datos de tutor y código de recuperación con su
		vigencia \\ \hline
	
		\texttt{quejas} & \texttt{queja-service} y \texttt{revision-\allowbreak service} & Cuerpo del expediente:
		folio, identidad del quejoso, narrativa de los hechos, unidad académica, denunciado, origen,
		estatus, y los campos del flujo de validación ---motivo de rechazo, área, defensor asignado,
		quién validó y cuándo--- \\ \hline
	
		\texttt{queja\_tutores} & \texttt{queja-service} & Datos del tutor o adulto responsable,
		en relación uno a uno con la queja \\ \hline
	
		\texttt{queja\_evidencias} & \texttt{queja-service} & Archivos de evidencia: nombre, tipo,
		tamaño y contenido binario completo \\ \hline
	
		\texttt{acuerdos\_} \texttt{conciliacion} & \texttt{revision-\allowbreak service} y
		\texttt{queja-service} & Acuerdos emitidos por la Defensoría: asunto, términos, estado y
		respuesta del quejoso \\ \hline
	
		\texttt{notificaciones} & \texttt{notificaciones-} \texttt{service} & Avisos persistidos por
		destinatario, con tipo, estado de lectura y enlace \\ \hline
	
		\texttt{dependencias} & \texttt{catalogo-\allowbreak service} & Catálogo institucional jerárquico con
		clave propia, clave del padre, tipo, categoría y nivel \\ \hline
	
		\texttt{personal\_} \texttt{administrativo} & \texttt{admin-service} & Cuentas de personal:
		identidad, rol, credencial cifrada, condición de cuenta temporal y último acceso \\ \hline
	
		\texttt{plantillas\_} \texttt{documentos} & \texttt{admin-service} & Plantillas de oficios
		con clave estable de negocio, contenido y datos de última actualización \\ \hline
	
		\texttt{bitacora\_acciones} & \texttt{admin-service} & Registro de auditoría: usuario, acción,
		dirección de origen y fecha \\ \hline
	
		\texttt{preguntas\_chatbot} & \texttt{chatbot-\allowbreak service} & Contenido del menú de preguntas
		frecuentes del portal público \\ \hline
	\end{longtable}
	
	\normalsize
	
	\subsubsection{Correspondencia entre reglas de negocio y mecanismos de integridad}
	Un modelo de datos que solo almacena información deja las reglas del dominio a merced de que
	cada consumidor las respete. El diseño de esta base busca lo contrario: que las reglas se hagan
	cumplir en el nivel más bajo posible, de modo que ningún camino de acceso pueda violarlas. La
	Tabla~\ref{tab:reglas_mecanismos} muestra esa correspondencia para las reglas que admiten
	verificación estructural.
	
	\small
	\renewcommand{\arraystretch}{1.2}
	
	\begin{longtable}{|p{1.9cm}|p{5.6cm}|p{6.6cm}|}
		\caption{Reglas de negocio y mecanismo que las hace cumplir}
		\label{tab:reglas_mecanismos} \\
		\hline
		\rowcolor{headerblue}
		\headercell{Regla} & \headercell{Enunciado} & \headercell{Mecanismo} \\
		\hline
		\endfirsthead
	
		\hline
		\rowcolor{headerblue}
		\headercell{Regla} & \headercell{Enunciado} & \headercell{Mecanismo} \\
		\hline
		\endhead
	
		RN-Q04 & El par folio-correo es la llave de activación & Consulta entre servicios que exige
		coincidencia de ambos campos antes de crear la cuenta \\ \hline
	
		RN-Q05 & Una cuenta activa no se reactiva & Verificación del indicador de actividad antes de
		procesar la activación \\ \hline
	
		RN-Q07 & El código de recuperación expira en diez minutos & Columna de fecha de expiración
		evaluada en el servidor al validar el código \\ \hline
	
		RN-Q08 & Una queja solo es editable en estatus recibida & Validación del estatus en la capa
		de negocio antes de aplicar la modificación \\ \hline
	
		RN-Q09 & Solo el destinatario responde un acuerdo & Comparación del correo del acuerdo contra
		el del token de sesión \\ \hline
	
		RN-Q10 & El origen de la queja se infiere & Asignación en el servidor según el punto de acceso
		invocado, sin exponerlo como campo de entrada \\ \hline
	
		RN-R01 & Transiciones válidas del expediente & Columna de estatus con conjunto acotado de
		valores y validación de la transición en la capa de negocio \\ \hline
	
		RN-R02 & Solo recepción usa el servicio de revisión & Declaración de rol requerido en cada
		controlador \\ \hline
	
		RN-R06 & Solo defensores activos reciben turnado & Consulta filtrada por rol y por indicador
		de actividad \\ \hline
	
		RN-A03 & La baja desactiva, no elimina & Indicador de actividad en lugar de borrado físico,
		preservando la integridad referencial con la bitácora \\ \hline
	
		RN-A04 & Correo y número de empleado únicos & Restricción de unicidad en la definición de la
		tabla \\ \hline
	
		RN-A06 & La plantilla se identifica por clave de negocio & Restricción de unicidad sobre la
		columna de tipo, independiente del identificador autogenerado \\ \hline
	\end{longtable}
	
	\normalsize
	\setlength{\tabcolsep}{6pt}
	\renewcommand{\arraystretch}{1}
	
	El diccionario de datos completo, con el tipo y la longitud de cada columna, las cardinalidades
	detalladas, los índices y las sentencias de creación, se documenta en el anexo correspondiente.
	
	\subsection{Diseño del clasificador de competencia}
	\label{sec:diseno_clasificador}
	El clasificador de competencia se diseñó como componente independiente de la plataforma, con la
	función de sugerir al personal de primer contacto si una queja es competencia de la Defensoría a
	partir de la narrativa de los hechos. Sus fundamentos conceptuales y su diseño detallado
	---la representación híbrida de características, el factor de amplificación de la señal
	normativa, la arquitectura del flujo de transformaciones y el protocolo de evaluación--- se
	desarrollan en las secciones~\ref{sec:problema_clasificacion} a~\ref{sec:validacion_robusta} del
	marco teórico, por lo que aquí se documenta únicamente su relación con el resto del sistema.
	
	El diseño de la integración prevé que el clasificador se consuma como un servicio adicional
	invocado por el flujo de validación, de modo que su ausencia no interrumpa la operación
	administrativa. Esta condición se recogió como requerimiento explícito: la indisponibilidad del
	componente inteligente no debe impedir que una queja se reciba, se valide o se turne. La
	consecuencia de diseño es que la sugerencia del clasificador es un dato adicional del
	expediente, nunca una precondición de su avance, lo cual además satisface la restricción de que
	la decisión final sea siempre humana.
	
	Al cierre de esta etapa el clasificador se desarrolló y evaluó de forma independiente ---los
	resultados experimentales se presentan en el Capítulo~\ref{cap:implementacion}--- y no forma
	parte de los siete microservicios desplegados. Su integración como servicio de la plataforma
	está prevista para la etapa siguiente.
	
	\subsection{Diseño de la interfaz de usuario}
	\label{sec:diseno_interfaz}
	El diseño de las interfaces partió de prototipos de alta fidelidad que se usaron como
	herramienta de validación temprana con el cliente. La razón de invertir en ellos antes de
	escribir código fue económica: corregir el orden de los campos de un formulario en un prototipo
	cuesta minutos, y hacerlo sobre una interfaz ya implementada y desplegada cuesta un ciclo
	completo de compilación y despliegue.
	
	Esa inversión se justificó en la práctica. La retroalimentación de las sesiones de revisión
	derivó en dos tipos de cambio que se incorporaron antes de implementar: ajustes en el orden y
	agrupación de los campos del formulario de queja, y la incorporación de mensajes de validación
	explícitos donde el prototipo original solo mostraba el campo en rojo.
	
	Los criterios de diseño que se fijaron para las tres aplicaciones fueron los siguientes:
	
	\begin{itemize}
		\item \textbf{Identidad institucional.} La paleta y la estructura del encabezado y del pie
		siguen la identidad visual del Instituto, de modo que el usuario reconozca que se encuentra en
		un servicio oficial y no en un sitio de terceros.
	
		\item \textbf{Formularios por pasos.} Los formularios extensos se dividen en secciones
		numeradas en lugar de presentarse como una lista continua de campos, para reducir la carga
		cognitiva y hacer visible el progreso.
	
		\item \textbf{Ayuda contextual no intrusiva.} Los avisos normativos que el usuario debe
		conocer ---como el plazo para presentar una queja--- se presentan en un elemento flotante
		consultable en cualquier momento, en lugar de un bloque fijo que reduce el espacio útil del
		formulario.
	
		\item \textbf{Retroalimentación explícita.} Toda operación produce un mensaje visible de
		éxito, error o validación pendiente. Este criterio se elevó a requerimiento no funcional
		(RNF-Q07) porque su ausencia fue la causa de que, en una versión temprana, los mensajes de
		error existieran en el código pero no fueran visibles para el usuario.
	
		\item \textbf{Consistencia entre paneles de personal.} La consola de administración y el panel
		de revisión comparten componentes, patrones de error y estilos base, para que el personal que
		use ambos no tenga que aprender dos lenguajes visuales distintos.
	\end{itemize}
	
	Los prototipos completos y las capturas de las interfaces implementadas se presentan en el
	Capítulo~\ref{cap:modulos} y en el anexo correspondiente.
